% Generated draft: requires mathematical and source review.
\documentclass[11pt]{article}
\usepackage[margin=25mm]{geometry}
\usepackage{fontspec}
\setmainfont{Latin Modern Roman}
\usepackage{amsmath,amssymb,mathtools,unicode-math}
\setmathfont{Latin Modern Math}
\usepackage{enumitem,needspace,xurl,hyperref}
\setlength{\parindent}{0pt}
\setlength{\parskip}{5pt}
\setlength{\emergencystretch}{3em}
\newcommand{\N}{\mathbb{N}}
\newcommand{\Z}{\mathbb{Z}}
\newcommand{\Q}{\mathbb{Q}}
\newcommand{\R}{\mathbb{R}}
\newcommand{\C}{\mathbb{C}}
\newcommand{\F}{\mathbb{F}}
\newcommand{\A}{\mathbb{A}}
\newcommand{\PP}{\mathbb{P}}
\newcommand{\E}{\mathbb{E}}
\newcommand{\eps}{\varepsilon}
\newcommand{\dd}{\,\mathrm{d}}
\begin{document}
\section*{UnsolvedMath ID: 20000969}
\textbf{Source problem code:} catalogue entry
\subsection{Definitions and mathematical statement}
Sam Hopkins's source question starts from a Venkatesh--Viswanath result, reference [34], relating maximal graph-parking functions to the full-support root space of the Kac--Moody algebra with Cartan matrix $2I-A(G)$. Hopkins asks whether all parking functions occur somewhere and whether the sandpile group acts on such a root space, and notes a natural maximal-parking-function basis labeling. Anton Dochtermann notes matching counts of last syzygies of the toppling ideal and chambers of a slice of the graphic arrangement, asking for structural links. The record supplies no full description of that arrangement slice.

For a precise version of the aid's formula, let $G$ be a connected finite simple graph with $n$ vertices and choose a sink $q$. A parking function is $f:V\setminus\{q\}\to\Z_{\geq0}$ such that every nonempty $U\subseteq V\setminus\{q\}$ has a vertex $v$ with $f(v)<\operatorname{outdeg}_U(v)$. Maximal means maximal coordinatewise among parking functions. In $S_0=\C[x_v:v\neq q]$, the parking ideal is
\[
 M_{G,q}=\left\langle\prod_{v\in U}x_v^{\operatorname{outdeg}_U(v)}:
             \varnothing\neq U\subseteq V\setminus\{q\}\right\rangle.
\]
Its standard monomials are exactly $x^f$ for all parking functions. For the homogeneous toppling ideal, put $S=\C[x_v:v\in V]$ and
\[
 I_G=\langle x^{u_+}-x^{u_-}:u\in\operatorname{im}_{\Z}\Delta_G\rangle,
\]
where $u=u_+-u_-$ and $\Delta_G$ is the graph Laplacian. Extending $M_{G,q}$ to $S$ preserves its total Betti numbers, but the finite standard-monomial basis statement concerns $S_0/M_{G,q}$, not an added free sink variable.

Let $\Pi$ run over partitions into $r$ nonempty blocks, each inducing a connected subgraph. Contract the blocks, discard loops and replace parallel edges by a single edge to obtain the simple support graph $Q_\Pi$. Let $\mathfrak g(2I-A(Q_\Pi))$ be its Kac--Moody algebra and $\delta_\Pi$ the sum of its $r$ simple roots. With total Betti numbers, the aid reports
\[
 \beta_{r-1}^{S_0}(S_0/M_{G,q})=\beta_{r-1}^{S}(S/I_G)
 =\sum_{\substack{\Pi\ \mathrm{connected}\ r\text{-}\mathrm{partition}}}
    \dim\mathfrak g(2I-A(Q_\Pi))_{\delta_\Pi},\qquad1\leq r\leq n.
\]
These are total, not multigraded, equalities. At $r=n$, the partition into singletons recovers the last-syzygy/root-space coincidence for $G$ itself. At other degrees the formula uses root spaces for connected quotient partitions, not one single root space containing every parking function.

The full parking-function set is also a torsor for the sandpile group
\[
 K(G,q)=\Z^{n-1}/\Delta_{G,q}\Z^{n-1},
\]
using its usual reduced-divisor/parking representative action. The aid states that this transitive action preserves the subset of maximal parking functions exactly for trees. For a connected graph with a cycle that subset is proper, so this usual action cannot restrict to the maximal-parking-function root basis. This does not rule out other representations of the sandpile group on a vector space; in particular it is not a nonexistence theorem for every possible root-space action. The all-degree numerical formula and this obstruction leave the broader requested structural Lie/algebra/arrangement connections open.
\subsection{Short English statement}
The aid extends the last-syzygy count to every total Betti number by summing root multiplicities over connected quotient partitions. All parking functions have a standard-monomial and torsor interpretation. The usual sandpile action, however, does not preserve the maximal-parking subset for a graph with a cycle, so it cannot simply act on that root basis by restriction.
\subsection{Sources}
\begin{itemize}
\item Local UnsolvedMath registry, permanent ID $20000969$, source code \texttt{catalogue entry}, supplied complete source record and research summary.
\item AIM workshop \emph{Generalizations of chip-firing and the critical group}, Problem (4), Sam Hopkins, with comments (a) by Hopkins and (b) by Anton Dochtermann; supplied source \url{https://aimath.org/pastworkshops/chipfiringproblems.pdf}. Canonical local record \texttt{aim-combinatorics-notes.json}, \texttt{notes[93]}. The supplied verification refers to PDF pages 2--3; no PDF was retrieved here.
\item Venkatesh and Viswanath [34], called a recent preprint by the historical source, without a supplied title, year or full entry. No full references for the all-degree Betti identity or arrangement-slice comparison are supplied.
\item The extraction also inserts a preceding-problem Eulerian-number footnote, \texttt{en.wikipedia.org/wiki/Eulerian number}, and a page header. The supplied verification identifies these as unrelated artifacts, not supporting references for this question.
\item Research attempt 1: partial result; problem partially solved at the recorded run; supplied rigor label ``proved'', novelty confidence ``low''. The machine-generated aid is not peer reviewed and requires independent verification before citation.
\end{itemize}
\end{document}
